\section{Finite holonomy, the Schläfli configuration, and \texorpdfstring{\(E_6\)}{E6} symmetry}
\label{p12-83-c20-s6:chap:holonomia-schlafli-e6-nuevo}
\index{finite holonomy}
\index{Schläfli}
\index{E6@\(E_6\)}

The Witt branch does not exhaust the exceptional structures accessible from
APP. The parallel branch of this chapter starts from the two mixed
polarizations sum--product and product--sum, together with a declared
plaquette orientation and central section. The computed curvature determines
the alternating form; this defines the central extension. An extension
is not inferred from cardinality 27 alone.

\subsection{Mixed APP curvatures}

On \(T_9=(\mathbb Z/9\mathbb Z)^2\), with
\(S(x,y)=x+y\) and \(P(x,y)=xy\), let
\[
 \Delta_xf=f(x+1,y)-f(x,y),\qquad
 \Delta_yf=f(x,y+1)-f(x,y).
\]
The mixed connections are
\[
 A^{SP}=(\Delta_xS,\Delta_yP)=(1,x),\qquad
 A^{PS}=(\Delta_xP,\Delta_yS)=(y,1).
\]
Their discrete curvature is
\[
 F_A=\Delta_xA_y-\Delta_yA_x.
\]

\begin{proposition}[Opposite oriented curvatures]
\label{p12-83-c20-s6:prop:curvaturas-app-e6-nuevo}
\[
 F_{A^{SP}}=1,\qquad F_{A^{PS}}=-1\pmod9.
\]
The circulation around an oriented rectangle with sides \(r,s\) is,
respectively, \(rs\) and \(-rs\).
\end{proposition}

\begin{proof}
\(\Delta_xx-\Delta_y1=1\) and
\(\Delta_x1-\Delta_yy=-1\). A rectangle contains \(rs\) plaquettes; on
summing, the interior edges cancel, leaving the boundary circulation.
\end{proof}

For displacements \(u=(a,b)\), \(v=(c,d)\), the alternating class is
\begin{equation}
 \sigma(u,v)=ad-bc\pmod9.
 \label{p12-83-c20-s6:eq:forma-alternante-e6-nueva}
\end{equation}
The orientation fixes the sign of \(\sigma\). The central section used
below belongs to the gauge of this realization; changing it by a coboundary
produces an isomorphic extension, not a new invariant.

\subsection{Central extension and ternary reduction}

We define the nonadic Heisenberg group
\begin{equation}
 \mathcal H_9=\{(a,b,z):a,b,z\in\mathbb Z/9\mathbb Z\},
 \quad
 (a,b,z)(c,d,w)=(a+c,b+d,z+w+bc).
 \label{p12-83-c20-s6:eq:heisenberg-9-nuevo}
\end{equation}
The commutator of two horizontal displacements is
\[
 [(a,b,0),(c,d,0)]=(0,0,bc-ad),
\]
which records \(-\sigma(u,v)\). Reducing each coordinate modulo three
yields
\begin{equation}
 H_3=\{(a,b,z):a,b,z\in\mathbb F_3\},\qquad |H_3|=27,
 \label{p12-83-c20-s6:eq:heisenberg-3-nuevo}
\end{equation}
with the same law. The kernel of
\(\mathcal H_9\to H_3\) is
\((3\mathbb Z/9\mathbb Z)^3\), of cardinality 27.

\subsection{Weyl--Heisenberg representation and central phase}

The nonadic extension admits an explicit unitary realization. Let
\[
 \zeta_9:=\exp(2\pi i/9),
 \qquad
 \mathscr H_9^{\rm Sch}:=\mathbb C^9,
\]
with basis \(\{|n\rangle:n\in\mathbb Z/9\mathbb Z\}\). We define
\begin{equation}
 X|n\rangle:=|n+1\rangle,
 \qquad
 Z|n\rangle:=\zeta_9^n|n\rangle.
 \label{p12-83-c20-s6:eq:weyl-desplazamiento-fase-u029}
\end{equation}
One has
\[
 ZX=\zeta_9XZ
\]
and, by multiplication,
\begin{equation}
 (X^aZ^b)(X^cZ^d)
 =\zeta_9^{bc}X^{a+c}Z^{b+d}.
 \label{p12-83-c20-s6:eq:cociclo-weyl-nonadico-u029}
\end{equation}

\begin{theorem}[Faithful Weyl--Heisenberg representation]
\label{p12-83-c20-s6:thm:representacion-weyl-heisenberg-u029}
The map
\begin{equation}
 \rho_{\rm WH}:\mathcal H_9\longrightarrow
 U(\mathscr H_9^{\rm Sch}),
 \qquad
 \rho_{\rm WH}(a,b,z):=\zeta_9^zX^aZ^b,
 \label{p12-83-c20-s6:eq:representacion-weyl-heisenberg-u029}
\end{equation}
is a faithful unitary representation. For
\(u=(a,b)\) and \(v=(c,d)\), its commutator satisfies
\begin{equation}
 [\rho_{\rm WH}(u,0),\rho_{\rm WH}(v,0)]
 =\zeta_9^{\,bc-ad}I
 =\zeta_9^{-\sigma(u,v)}I.
 \label{p12-83-c20-s6:eq:conmutador-weyl-holonomia-u029}
\end{equation}
\end{theorem}

\begin{proof}
Identity \eqref{p12-83-c20-s6:eq:cociclo-weyl-nonadico-u029} reproduces the term
\(bc\) of \eqref{p12-83-c20-s6:eq:heisenberg-9-nuevo}; hence the homomorphism
property. If \(\zeta_9^zX^aZ^b=I\), the action on the basis successively
forces \(a=0\), \(b=0\), and \(z=0\). This proves faithfulness. Comparing
the products in both orders gives
\(\zeta_9^{bc-ad}\), which is the alternating phase of
\eqref{p12-83-c20-s6:eq:forma-alternante-e6-nueva}.
\end{proof}

The representation distinguishes two conjugate observables. Let
\[
 Q_{\rm H}:=\operatorname{diag}(0,1,\ldots,8),
 \qquad
 F_{jk}:=\frac13\zeta_9^{jk},
 \qquad
 P_{\rm H}:=F^\dagger Q_{\rm H}F.
\]
The eigenbases of \(Q_{\rm H}\) and \(P_{\rm H}\) are mutually
unbiased, since \(|F_{jk}|^2=1/9\). For every unit vector
\(\psi\in\mathscr H_9^{\rm Sch}\),
\begin{align}
 \operatorname{Var}_\psi(Q_{\rm H})
 \operatorname{Var}_\psi(P_{\rm H})
 &\ge
 \frac14\left|
 \langle\psi,[Q_{\rm H},P_{\rm H}]\psi\rangle
 \right|^2,
 \label{p12-83-c20-s6:eq:incertidumbre-weyl-u029}\\
 H_{Q_{\rm H}}(\psi)+H_{P_{\rm H}}(\psi)
 &\ge\log9.
 \label{p12-83-c20-s6:eq:incertidumbre-entropica-weyl-u029}
\end{align}
Here \(H_{Q_{\rm H}}\) and \(H_{P_{\rm H}}\) are the Shannon entropies
of the Born distributions in the respective eigenbases,
with natural logarithm and convention \(0\log0=0\).
These operators are dimensionless. A physical units chart is a
subsequent realization and does not intervene in the central extension.

\subsection{Action of 10 displacements on the Cayley graph}

For a linear functional
\(\ell(a,b)=pa+qb\), set
\[
 D_\ell(a,b)=\left(a,b,\frac{ab}{2}+\ell(a,b)\right),
 \qquad Z=(0,0,1),
\]
where \(2^{-1}=2\) in \(\mathbb F_3\). The inverse-closed set
\begin{equation}
 S_\ell=
 \{D_\ell(v):v\in\mathbb F_3^2\setminus\{0\}\}
 \cup\{Z,Z^{-1}\}
 \label{p12-83-c20-s6:eq:conjunto-conexion-e6}
\end{equation}
has ten elements and does not contain the identity.

\begin{theorem}[Independence of the linear gauge]
\label{p12-83-c20-s6:thm:independencia-calibre-e6}
The nine functionals \(\ell\) produce isomorphic Cayley graphs. The
nine sets \(S_\ell\) form a single orbit under
\(\operatorname{Aut}(H_3)\).
\end{theorem}

\begin{proof}
The map
\[
 \phi_\ell(a,b,z)=(a,b,z+\ell(a,b))
\]
is a central automorphism and satisfies
\(\phi_\ell(S_0)=S_\ell\). To write the complete collection in the
polarized coordinates of \cref{p12-83-c20-s6:eq:heisenberg-9-nuevo}, let
\(c(v,w)=b c\) for \(v=(a,b)\), \(w=(c,d)\). Given
\(M\in\operatorname{GL}(2,3)\), the defect
\[
 B_M(v,w)=c(Mv,Mw)-\det(M)c(v,w)
\]
is symmetric, because the alternating parts of the two terms coincide.
We define the quadratic correction
\begin{equation}
 q_M(v)=2B_M(v,v),
 \qquad
 q_M(v+w)-q_M(v)-q_M(w)=B_M(v,w),
 \label{p12-83-c20-s6:eq:correccion-cuadratica-automorfismos-h3}
\end{equation}
where \(2=2^{-1}\) in \(\mathbb F_3\). The complete collection of
automorphisms is then written
\begin{equation}
 (v,t)\longmapsto
 \bigl(Mv,\det(M)t+q_M(v)+\ell(v)\bigr),
 \qquad M\in\operatorname{GL}(2,3),\
 \ell\in(\mathbb F_3^2)^*.
 \label{p12-83-c20-s6:eq:automorfismos-h3-completos}
\end{equation}
The polar identity of \cref{p12-83-c20-s6:eq:correccion-cuadratica-automorfismos-h3}
directly proves that the central product is preserved. The collection has
\(48\cdot9=432\) elements. Exhaustive enumeration of the
\(\binom{13}{5}=1287\) inverse-closed subsets of cardinality ten
finds exactly those nine with the parameters derived
below.
\end{proof}

Write \(\underline S=\sum_{s\in S}s\) in
\(\mathbb Z[H_3]\). Counting ordered products gives
\begin{equation}
 \underline S^{\,2}
 =10e+\underline S+5(\underline{H_3}-e-\underline S)
 =5\underline{H_3}+5e-4\underline S.
 \label{p12-83-c20-s6:eq:identidad-grupo-schlafli-nueva}
\end{equation}
Indeed, the identity receives ten factorizations; an element of
\(S\) receives one; each of the other sixteen receives five.

\begin{theorem}[Graph of the twenty-seven lines]
\label{p12-83-c20-s6:thm:grafo-27-rectas-nuevo}
The adjacency matrix \(A_\ell\) of
\(\Gamma_\ell=\operatorname{Cay}(H_3,S_\ell)\) satisfies
\begin{equation}
 A_\ell^2=5I-4A_\ell+5J.
 \label{p12-83-c20-s6:eq:identidad-schlafli-nueva}
\end{equation}
Therefore,
\[
 \Gamma_\ell=\operatorname{srg}(27,10,1,5),
 \qquad
 \operatorname{Spec}(A_\ell)=\{10^1,1^{20},(-5)^6\}.
\]
It is isomorphic to the intersection graph of the twenty-seven lines on a
smooth cubic surface.
\end{theorem}

\begin{proof}
The regular representation of
\cref{p12-83-c20-s6:eq:identidad-grupo-schlafli-nueva} produces
\cref{p12-83-c20-s6:eq:identidad-schlafli-nueva}. Its entries say that every vertex
has ten neighbours, that an adjacent pair has one common neighbour, and a nonadjacent
pair has five. On the constant line, \(A_\ell\) acts by 10;
on its complement, \(J=0\) and
\((A_\ell-I)(A_\ell+5I)=0\). Zero trace fixes the multiplicities 20 and 6.

In the model of a cubic obtained by blowing up six points of
\(\mathbb P^2\), the 27 classes are
\[
 E_i,\qquad L_{ij}=H-E_i-E_j,\qquad
 Q_i=2H-\sum_{j\ne i}E_j.
\]
With \(H^2=1\), \(E_i^2=-1\), \(H\cdot E_i=0\), and
\(E_i\cdot E_j=0\) for \(i\ne j\), the adjacency matrix defined
by intersection between distinct classes has the same parameters.
The complete intersection matrix has diagonal \(-1\).
Identification with the classical configuration uses the
uniqueness theorem for this realization of the 27 lines
\cite{Manivel2005}, with the domain and action specified there.
\end{proof}

The isomorphism statement can be made fully effective through
a computed bijection. The labeling is not absolute: another linear section
or an automorphism of the surface produces another table in the same class.
\begin{longtable}{@{}c@{\,}c@{\,}c@{\qquad}l@{}}
\caption{Explicit bijection between \(H_3\) and the twenty-seven line
classes.}
\label{p12-83-c20-s6:tab:h3-27-rectas-u029}\\
\toprule
\(a\)&\(b\)&\(z\)&line class\\
\midrule
\endfirsthead
\multicolumn{4}{@{}l}{\small\itshape Table \thetable\ (continued)}\\
\toprule
\(a\)&\(b\)&\(z\)&line class\\
\midrule
\endhead
0&0&0&\(E_1\)\\
0&0&1&\(L_{12}\)\\
0&0&2&\(Q_2\)\\
0&1&0&\(L_{13}\)\\
0&1&1&\(Q_1\)\\
0&1&2&\(E_3\)\\
0&2&0&\(Q_3\)\\
0&2&1&\(E_2\)\\
0&2&2&\(L_{23}\)\\
1&0&0&\(L_{14}\)\\
1&0&1&\(L_{35}\)\\
1&0&2&\(L_{26}\)\\
1&1&0&\(L_{36}\)\\
1&1&1&\(L_{24}\)\\
1&1&2&\(L_{15}\)\\
1&2&0&\(L_{25}\)\\
1&2&1&\(L_{16}\)\\
1&2&2&\(L_{34}\)\\
2&0&0&\(Q_4\)\\
2&0&1&\(L_{46}\)\\
2&0&2&\(E_6\)\\
2&1&0&\(E_5\)\\
2&1&1&\(Q_6\)\\
2&1&2&\(L_{56}\)\\
2&2&0&\(L_{45}\)\\
2&2&1&\(E_4\)\\
2&2&2&\(Q_5\)\\
\bottomrule
\end{longtable}

As a local check, the ten neighbours of \(e=(0,0,0)\) are the elements of
\(S_0\). The table transforms them into the ten classes intersecting \(E_1\):
the five \(L_{1j}\) and the five \(Q_j\), with \(2\le j\le6\). The
complete verification compares all 729 entries of both adjacency
matrices.

The complement \(B_\ell=J-I-A_\ell\) satisfies
\begin{equation}
 B_\ell^2=8I+2B_\ell+8J,
 \qquad
 \operatorname{srg}(27,16,10,8),
 \qquad
 \operatorname{Spec}(B_\ell)=\{16^1,4^6,(-2)^{20}\}.
 \label{p12-83-c20-s6:eq:complemento-schlafli}
\end{equation}
This is the Schläfli graph under the complementary adjacency convention;
\(A_\ell\) is the intersection graph with parameters
\((27,10,1,5)\).

Each edge belongs to a unique triangle, since \(\lambda=1\). The number
of triangles is
\[
 \frac{27\cdot10}{6}=45.
\]
The full automorphism group has order
\[
 27\cdot1920=51840
\]
and, through its action on the line configuration, is identified with
\(W(E_6)\), not merely through equality of orders \cite{Manivel2005}.

\subsection{The root lattice realizing \texorpdfstring{\(E_6\)}{E6}}

To make explicit the object on which that group acts, let
\begin{equation}
 \operatorname{Pic}(X)
 =\mathbb ZH\oplus\bigoplus_{i=1}^6\mathbb ZE_i,
 \qquad
 H^2=1,\qquad E_i^2=-1,\qquad H\cdot E_i=E_i\cdot E_j=0\ (i\ne j),
 \label{p12-83-c20-s6:eq:picard-cubica-e6}
\end{equation}
be the Picard lattice of the blow-up of six general points of
\(\mathbb P^2\). Its canonical class is
\[
 K_X=-3H+E_1+\cdots+E_6.
\]
The orthogonal complement
\begin{equation}
 L_{E_6}=K_X^\perp\subset\operatorname{Pic}(X),
 \qquad (x,y)_{E_6}:=-x\cdot y,
 \label{p12-83-c20-s6:eq:reticulo-raices-e6}
\end{equation}
is positive definite of rank six. A basis of simple roots is
\begin{equation}
 \begin{aligned}
 \alpha_1&=E_1-E_2,& \alpha_2&=E_2-E_3,&
 \alpha_3&=E_3-E_4,\\
 \alpha_4&=E_4-E_5,& \alpha_5&=E_5-E_6,&
 \alpha_6&=H-E_1-E_2-E_3.
 \end{aligned}
 \label{p12-83-c20-s6:eq:raices-simples-e6}
\end{equation}
Each \(\alpha_i\) is orthogonal to \(K_X\), has norm two, and their products
are \(-1\) exactly for the adjacent pairs in the diagram
\(E_6\): the chain \(1-2-3-4-5\), with node \(6\) attached to node \(3\).
Therefore, its Gram matrix is the Cartan matrix of \(E_6\).

\begin{proposition}[The seventy-two roots and their reflections]
\label{p12-83-c20-s6:prop:setenta-dos-raices-e6}
The roots of \(L_{E_6}\) are
\begin{equation}
 \begin{aligned}
 &E_i-E_j &&(i\ne j),\\
 &\pm(H-E_i-E_j-E_k) &&(1\le i<j<k\le6),\\
 &\pm(2H-E_1-\cdots-E_6),
 \end{aligned}
 \label{p12-83-c20-s6:eq:enumeracion-raices-e6}
\end{equation}
a total of \(30+40+2=72\). For each root \(\alpha\),
\begin{equation}
 s_\alpha(x)=x+(x\cdot\alpha)\alpha
 \label{p12-83-c20-s6:eq:reflexion-raiz-e6}
\end{equation}
preserves intersection and \(K_X\). The six simple reflections generate
\(W(E_6)\) and faithfully permute the 27 classes
\(E_i,L_{ij},Q_i\) of \cref{p12-83-c20-s6:thm:grafo-27-rectas-nuevo}.
\end{proposition}

\begin{proof}
Substituting each class from \cref{p12-83-c20-s6:eq:enumeracion-raices-e6} into
\cref{p12-83-c20-s6:eq:picard-cubica-e6} gives \(K_X\cdot\alpha=0\) and
\(\alpha^2=-2\). The count of classes is as stated. The reflection formula
is the orthogonal reflection for the form \(-(\cdot)\); it preserves \(K_X\) because
\(K_X\cdot\alpha=0\). The simple system in
\cref{p12-83-c20-s6:eq:raices-simples-e6} generates all the roots, and its Coxeter group is
\(W(E_6)\). Its action on the 27 classes of exceptional curves preserves
intersection, so it coincides with the action on the graph already
constructed \cite{Manivel2005}.
\end{proof}

\subsection{Stratification of the 729 words}

We write a six-trit word as
\[
 w=(\ell_1,h_1,\ell_2,h_2,\ell_3,h_3),
\]
and define two elements of \(H_3\),
\[
 L(w)=(\ell_1,\ell_2,\ell_3),\qquad
 H(w)=(h_1,h_2,h_3),
\]
using the central law in their products. According to
\(L(w)^{-1}H(w)\), the 729 words are divided into
\begin{align*}
 \mathcal D_\ell&=\{w:L(w)^{-1}H(w)=e\},\\
 \mathcal I_\ell&=\{w:L(w)^{-1}H(w)\in S_\ell\},\\
 \mathcal S_\ell&=\{w:L(w)^{-1}H(w)\notin\{e\}\cup S_\ell\}.
\end{align*}

\begin{theorem}[Hensel--Schläfli partition]
\label{p12-83-c20-s6:thm:particion-hensel-schlafli-nueva}
\begin{equation}
 729=27+270+432,
 \label{p12-83-c20-s6:eq:particion-729-e6-nueva}
\end{equation}
corresponding to the diagonal, incidence, and skewness.
\end{theorem}

\begin{proof}
For each of the 27 values of \(L(w)\), there is a unique value of
\(H(w)\) with identity difference, ten with difference in \(S_\ell\), and
sixteen remaining. Multiply \(27\) by \(1,10,16\).
\end{proof}

The cardinality \(432\) of the skew stratum coincides numerically with the
multiplicity of the transverse region of \(\pi\) in
\cref{pi-estrella:figura}. The two objects are not thereby identified: here
one counts words classified by \(L(w)^{-1}H(w)\), whereas there one
counts preimages of a register \(U_6\). Without an explicit equivariant
intertwiner between those domains, \(432=432\) is only a cardinality check.

% BEGIN REV09_RECTAS27_COORDENADAS
\subsection{Reversible coordinates for ordered pairs of lines}
\label{corpus9:xii:rectas-coordenadas}

Table~\ref{p12-83-c20-s6:tab:h3-27-rectas-u029} fixes the bijection
\(\iota:H_3\to\mathscr L_{27}\), where \(\mathscr L_{27}\) is the
configuration of twenty-seven lines. We retain the central law
\eqref{p12-83-c20-s6:eq:heisenberg-3-nuevo}:
\[
 (a,b,z)(c,d,w)=(a+c,b+d,z+w+bc),\qquad
 (a,b,z)^{-1}=(-a,-b,-z+ab).
\]
Multiplication in both orders verifies the inverse formula: all three
resulting coordinates vanish. All operations in this subsection are
performed in \(\mathbb F_3\).

\begin{proposition}[Reversible encoding of the 729 words]
\label{corpus9:xii:rectas-biyeccion}
For \(w=(\ell_1,h_1,\ell_2,h_2,\ell_3,h_3)\), let
\(L(w)=(\ell_1,\ell_2,\ell_3)\) and
\(H(w)=(h_1,h_2,h_3)\). The map
\begin{equation}
 \mathcal E:\mathbb F_3^6\longrightarrow
                 \mathscr L_{27}\times\mathscr L_{27},\qquad
 \mathcal E(w)=(\iota L(w),\iota H(w))
 \label{corpus9:xii:rectas-codificacion}
\end{equation}
is bijective. In the relative coordinates
\(L=(a,b,z)\), \(R=L^{-1}H=(u,v,t)\), reconstruction is given by
\begin{equation}
 H=(a+u,b+v,z+t+bu),\qquad
 w=(a,a+u,b,b+v,z,z+t+bu).
 \label{corpus9:xii:rectas-reconstruccion}
\end{equation}
\end{proposition}
\begin{proof}
Given an ordered pair \((\ell,h)\) of lines, the table fixes the
unique triples \(L=\iota^{-1}(\ell)\) and
\(H=\iota^{-1}(h)\). Interlacing their coordinates recovers a unique
word \(w\), and evaluation of
\eqref{corpus9:xii:rectas-codificacion} returns the original pair.
Conversely, separating and then interlacing a word preserves all six
coordinates. This establishes both inverse compositions.

The transformation \((L,H)\mapsto(L,L^{-1}H)\) is bijective, with
inverse \((L,R)\mapsto(L,LR)\). The central law gives
\(LR=(a+u,b+v,z+t+bu)\), proving
\eqref{corpus9:xii:rectas-reconstruccion}. In particular, the term
\(bu\) retains the cocycle in the third coordinate. Omitting it changes
the reconstruction by precisely \(bu\); for example,
\(L=(0,1,0)\) and \(R=(1,0,0)\) produce
\(H=(1,1,1)\), whereas componentwise addition produces
\((1,1,0)\).
\end{proof}

The partition in
Theorem~\ref{p12-83-c20-s6:thm:particion-hensel-schlafli-nueva}
is preserved by this bijection: the relative element \(R\) distinguishes
the diagonal, incident, and skew strata, of cardinalities \(27\),
\(270\), and \(432\). Ordered pairs constitute the full domain of
the chart. Adjacency is an additional condition on pairs or their
sequences. Prolongation to complete histories is constructed in
Theorem~\ref{corpus9:xii:rectas-homeomorfismo} through the TPK block
reader.
% END REV09_RECTAS27_COORDENADAS


\subsection{Necessity of the central phase}

Enumerating all inverse-closed subsets of cardinality ten
in the five groups of order 27 gives
\[
\begin{array}{c|c|c}
\text{group}&\text{abelian}&
\#\operatorname{srg}(27,10,1,5)\\ \hline
C_{27}&\text{yes}&0\\
C_9\times C_3&\text{yes}&0\\
C_3^3&\text{yes}&0\\
C_9\rtimes C_3&\text{no}&9\\
H_3&\text{no}&9.
\end{array}
\]
Thus, preserving 27 labels but abelianizing composition destroys the
configuration. Neither cardinality nor degree suffices.

If \(\mathcal T\) is the set of 45 triangles, the cubic
\[
 \mathcal C_\Gamma(x_1,\ldots,x_{27})
 =\sum_{\{i,j,k\}\in\mathcal T}x_ix_jx_k
\]
recovers adjacency, because each edge lies in a unique triangle. Its
permutation stabilizer is exactly
\[
 \operatorname{Aut}(\Gamma_\ell)\cong W(E_6).
\]

For the following diagram we write
\(\Gamma_{27}=\Gamma_\ell\) and \(\mathcal T_{45}=\mathcal T\).
\begin{figure}[htbp]
\centering
\[
 \begin{gathered}
 \text{APP curvatures }\pm1
 \longrightarrow\sigma
 \longrightarrow\mathcal H_9
 \longrightarrow H_3,\\[3pt]
 H_3
 \longrightarrow\Gamma_{27}
 \longrightarrow\mathcal T_{45}
 \longrightarrow W(E_6).
 \end{gathered}
\]
\caption{Causal chain of the mixed-holonomy branch. It is parallel to
\(C_W\to W_{12}\to M_{12}\), not a link within it.}
\label{p12-83-c20-s6:fig:cadena-holonomia-e6-u029}
\end{figure}

\begin{criterion}[Refutation criteria for the \(E_6\) branch]
\label{p12-83-c20-s6:crit:refutacion-rama-e6-u029}
The branch is refuted if the curvatures are not opposite, if the central
cocycle is erased, if the quadratic correction
\eqref{p12-83-c20-s6:eq:correccion-cuadratica-automorfismos-h3} does not preserve the product,
if \eqref{p12-83-c20-s6:eq:identidad-grupo-schlafli-nueva} does not produce the coefficients
\((10,1,5)\), if an abelian group produces the same configuration, if the
partition does not sum to 729, if Table
\ref{p12-83-c20-s6:tab:h3-27-rectas-u029} does not intertwine the adjacency matrices, or if
\(W(E_6)\) is identified solely by its order and not by its action on the
twenty-seven classes and the 72 roots.
\end{criterion}
